% Image credits for the photographs and AI illustrations of Book 3
% (University Chemistry, Year 2), French edition. Machine-readable license
% records live in images/book3/CREDITS.md; keep this file in sync with the
% English frontmatter/image-credits-book3.tex.
\chapter*{Crédits des images}

Les dessins de ce livre sont originaux. Les photographies sont utilisées
sous leurs licences libres respectives, avec nos remerciements à leurs
auteurs~; les liens vers les sources et les adresses des licences de
chaque photographie sont recensés dans
\texttt{images/book3/CREDITS.md}, dans le dépôt du livre.
Les pictogrammes de danger sont ceux du Système général harmonisé de
classification et d'étiquetage des produits chimiques, tels que les
dessine l'extension \texttt{ghsystem} de \TeX\ Live.

\bigskip
Outre les photographies et les dessins originaux, certaines figures sont
des illustrations générées par intelligence artificielle, produites avec
le générateur d'images d'OpenAI à partir de consignes rédigées par les
éditeurs du livre, et relues pour leur exactitude chimique avant
inclusion. La consigne complète de chaque image générée est consignée
dans \texttt{images/book3/ai/PROMPTS.md}, dans le dépôt.

\bigskip
\noindent Photographies, par chapitre~:
\begin{itemize}\small
  \item Chap.~1~: Germain Henri Hess, portrait par I.~I. Reimers
        (1837--1838), CC BY-SA 4.0~; Marcellin Berthelot, auteur inconnu,
        domaine public (tous deux sur Wikimedia Commons).
  \item Chap.~2~: Josiah Willard Gibbs, photographe inconnu, domaine
        public (Wikimedia Commons).
  \item Chap.~3~: François-Marie Raoult, photographe inconnu, domaine
        public (Wikimedia Commons).
  \item Chap.~4~: Jacobus Henricus van 't Hoff, par Nicola Perscheid
        (1904), domaine public~; Henry Le Chatelier, auteur inconnu,
        CC BY-SA 4.0~; Fritz Haber et Carl Bosch, Fondation Nobel,
        domaine public (tous sur Wikimedia Commons).
  \item Chap.~6~: haut fourneau de Duisburg-Nord, par Dietmar Rabich,
        CC BY-SA 4.0~; soudage aluminothermique d'un rail, par PetrS.,
        CC BY-SA 3.0 (tous deux sur Wikimedia Commons).
  \item Chap.~7~: colonne de distillation d'une raffinerie, par CEphoto,
        Uwe Aranas, CC BY-SA 3.0 (Wikimedia Commons).
  \item Chap.~9~: Michael Faraday, par Thomas Phillips (1842), domaine
        public (Wikimedia Commons).
  \item Chap.~10~: Jaroslav Heyrovský, archives de l'Institut de chimie
        physique J.~Heyrovský, CC BY-SA 3.0 (Wikimedia Commons).
  \item Chap.~11~: une pile voltaïque, par GuidoB, CC BY-SA 3.0~;
        Charles Martin Hall et Paul Héroult, auteurs inconnus, domaine
        public (tous sur Wikimedia Commons).
  \item Chap.~12~: une statue en cuivre, par Elcobbola, domaine public~;
        une anode sacrificielle, par Zwergelstern, CC BY-SA 3.0 (tous deux
        sur Wikimedia Commons).
  \item Chap.~13~: Erwin Schrödinger, Fondation Nobel (1933), domaine
        public (Wikimedia Commons).
  \item Chap.~14~: dioxygène liquide dans un aimant, par Pieter Kuiper,
        domaine public~; Robert Mulliken et Friedrich Hund, par GFHund,
        CC BY 3.0 (tous deux sur Wikimedia Commons).
  \item Chap.~16~: August Kekulé (1873), auteur inconnu, domaine public
        (Wikimedia Commons).
  \item Chap.~18~: solutions de métaux de transition, par Benjah-bmm27,
        domaine public~; Alfred Werner, Les Prix Nobel (1914), domaine
        public (tous deux sur Wikimedia Commons).
  \item Chap.~19~: un cristal de rubis, domaine public~; une émeraude, par
        Didier Descouens, CC BY-SA 3.0~; cristaux de sulfate de cuivre, par
        Stephanb, CC BY-SA 3.0 (tous sur Wikimedia Commons).
  \item Chap.~20~: Akira Suzuki, Ei-ichi Negishi et Richard Heck, par
        BloodIce, CC BY-SA 4.0 (Wikimedia Commons).
  \item Chap.~22~: Charles Friedel (années 1890) et James Mason Crafts
        (\textit{Popular Science Monthly}, 1900), domaine public (Wikimedia
        Commons).
  \item Chap.~24~: Michel Eugène Chevreul, par Nicolas Eustache Maurin,
        domaine public (Wikimedia Commons).
  \item Chap.~25~: Alexandre Borodine, auteur inconnu, domaine public~;
        Ludwig Claisen, par Veronica.villagrasa, CC BY-SA 4.0 (tous deux sur
        Wikimedia Commons).
  \item Chap.~26~: Robert Robinson, auteur inconnu, domaine public
        (Wikimedia Commons).
  \item Chap.~28~: Elias James Corey, par Trvthchem, CC0 (Wikimedia
        Commons).
  \item Chap.~29~: Wallace Carothers au laboratoire, photographe inconnu,
        domaine public (Wikimedia Commons).
  \item Chap.~30~: Emil Fischer, Atelier Victoria, Berlin, vers 1895,
        domaine public (Wikimedia Commons).
  \item Chap.~31~: Mikhaïl Tsvet, auteur inconnu, domaine public
        (Wikimedia Commons).
  \item Chap.~32~: Francis William Aston (1922), auteur inconnu, domaine
        public~; Gustav Kirchhoff, Robert Bunsen et Henry Roscoe (1862),
        domaine public (tous deux sur Wikimedia Commons).
  \item Chap.~34~: William Sealy Gosset (1908), domaine public (Wikimedia
        Commons).
\end{itemize}
\clearpage
